\documentclass[11pt,a4paper]{article}

\usepackage{amsmath,amssymb,amsthm}
\usepackage[colorlinks=true,urlcolor=blue,citecolor=blue]{hyperref}
\usepackage{geometry}
\usepackage{booktabs}
\geometry{margin=2.5cm}

\title{\textbf{Corrigendum}\\[0.5em]
	\large Theorem~2 of ``Markov-based continuous learning with diversion of data distribution
	direction for streaming data in limited memory''}

\author{Peemapat Wongsriphisant \and Kitiporn Plaimas \and Chidchanok Lursinsap}

\date{%
	\textit{Expert Systems With Applications}, 298 (2026) 129818\\
	DOI: \href{https://doi.org/10.1016/j.eswa.2025.129818}{10.1016/j.eswa.2025.129818}
}

\begin{document}
	\maketitle
	
	\noindent
	This notice addresses a mathematical error in Theorem~2 of the originally published
	paper and provides the corrected theoretical mechanism underlying the $D^4$ algorithm.
	\textbf{This theoretical correction does not affect the algorithm's implementation,
		functionality, or the experimental results reported in the paper.}
	
	% ──────────────────────────────────────────────────────────────────────────────
	\section{Correction to Theorem~2}
	
	In the original paper, the proof for Theorem~2 contains an algebraic sign error in
	the final steps. When taking the reciprocal and subtracting from a constant, the
	inequality sign must be reversed. Consequently, the original conclusion stating that
	the distance between the projected hyper-ellipsoids increases
	($d(\hat{\Omega}_a, \hat{\Omega}_b) \geq d(\Omega_a, \Omega_b)$) is incorrect.
	
	The mathematically corrected conclusion is the opposite: the projected distance is always less than or equal to the original distance,
	\begin{equation}
		d(\hat\Omega_a,\hat\Omega_b) \;\le\; d(\Omega_a,\Omega_b). \tag{10, corrected}
	\end{equation}
	Therefore, the resolution of classification ambiguity is not achieved by proving the hyper-ellipsoids become further apart in terms of Definition~2, but by substituting the original, unreliable classification scores with a new set of projected scores that exhibit greater separability.
	
	% ──────────────────────────────────────────────────────────────────────────────
	\section{Corrected Mechanism: Differential Penalty}
	
	The theoretical justification for $D^4$ relies on the asymmetric reduction of the metric scores. Let $\psi_a(\mathbf{x})$ be the original metric for a hyper-ellipsoid $\Omega_a$, and $\psi'_a(\mathbf{x})$ be the projected metric after selecting a subset of eigenvector indices $D_a$.

	The original metric can be partitioned into the projected metric and a non-negative \emph{removed penalty} term $\Delta_a(\mathbf{x})$ representing the summation over the discarded axes ($D'_a$):
	\[
	\Delta_a(\mathbf{x}) = \sum_{j \in D'_a} \frac{((\mathbf{x}-\mathbf{c}_a)^T \mathbf{u}_{a,j})^2}{w_{a,j}^2}
	\]
	Since $w_{a,j}^2 > 0$, it is guaranteed that $\Delta_a(\mathbf{x}) \ge 0$. The relationship between the scores simplifies to:
	\[
	\psi_a(\mathbf{x}) \;=\; \psi'_a(\mathbf{x}) \;+\; \Delta_a(\mathbf{x})
	\]
	It follows directly that the projected score is always less than or equal to the original score ($\psi'_a(\mathbf{x})\le\psi_a(\mathbf{x})$).

	The effectiveness of $D^4$ arises because this reduction is \textbf{not uniform} for both classes. In the overlap region (where $\psi_a < 0$ and $\psi_b < 0$), the algorithm substitutes the ambiguous decision $\arg\min(\psi_a,\psi_b)$ with a more robust decision:
	\[
	\arg\min\!\left(\psi'_a,\psi'_b\right)
	\;\equiv\;
	\arg\min\!\left(\psi_a-\Delta_a,\;\psi_b-\Delta_b\right)
	\]

	A central design goal of $D^4$ is to enhance the score separability, $|\psi'_a(\mathbf{x})-\psi'_b(\mathbf{x})|$. This is achieved when the \emph{differential penalty} $|\Delta_a(\mathbf{x})-\Delta_b(\mathbf{x})|$ amplifies the original separation.

	The $D^4$ \textbf{Compactness criterion} serves as the primary heuristic for achieving favorable differential penalties. By prioritizing the selection of eigenvectors corresponding to the smallest eigenvalues, it guarantees that the discarded axes ($D'_a$) are those with the largest widths ($w_{a,j}^2$). Maximizing the denominators in the penalty term systematically reduces the per-unit-projection contribution of each removed axis. This geometric criterion exploits the natural spread of the classes to create asymmetric penalties, transforming an ambiguous classification into a decisive one.
	
	% ──────────────────────────────────────────────────────────────────────────────
	\section{Impact Summary}
	
	\begin{table}[h]
		\centering
		\small
		\begin{tabular}{@{}l l p{5.5cm}@{}}
			\toprule
			\textbf{Component} & \textbf{Status} & \textbf{Reason} \\
			\midrule
			Algorithm implementation & Unaffected & Theorem~2 concerns analysis only \\
			Experimental results & Unaffected & No algorithmic change \\
			Theorem~1 (Parallelism) & Valid & Independent proof \\
			Theorem~3 (Closeness) & Valid & Uses triangle inequality only \\
			\textbf{Theorem~2 (Eq.~10)} & \textbf{Replaced} & \textbf{See Section~2} \\
			\bottomrule
		\end{tabular}
	\end{table}
	
	\noindent
	The corrected theoretical framework provides a more precise characterization of the $D^4$ mechanism. All conclusions drawn from the experimental evaluation in the original paper remain valid.
	
\end{document}